\chapter{T}

\section{TT}

    \subsection{Observation}\label{chVI-3-40}
        The slice constructions assemble into functors
        \begin{equation*}
            \ls, \rs \colon \TwAr(\sSet) \to \sSet \,.
        \end{equation*}
        In particular, for every sequence of map
        \begin{equation*}
            A \xrightarrow{i} X \xrightarrow{f} Y \xrightarrow{p} B
        \end{equation*}	
        we obtain maps
        \begin{align*}
            Y_{f /} \to Y_{f i /} \underset{B_{p f i /}}{\times} B_{p f /} \\
            Y_{/ f} \to Y_{/ f i} \underset{B_{/ p f i}}{\times} B_{p f} \,.
        \end{align*}
        It is these maps that appear in the proof of Joyal's theorem.
        \smallskip

        \emph{Proof}: The first part is immediate from the discussion preceding \ref{chVI-3-39}. For the second one we need to show that
        \begin{equation*}
            \begin{tikzcd}
                Y_{f/} \arrow[d, "p_{\star}"] \arrow[r, "i^{\star}"] & Y_{f i /} \arrow[d, "p_{\star}"] \\
                B_{p f /} \arrow[r, "i^{\star}"] & B_{p f i /}
            \end{tikzcd}	
        \end{equation*}
        commutes. But this is immediate from the compositions
        \begin{equation*}
            \begin{tikzcd}
                X \arrow[r, "f"] & Y \arrow[d, equal] \\
                A \arrow[d, equal] \arrow[u, "i"] \arrow[r, "f i"] & Y \arrow[d, "p"] \\
                A \arrow[r, "p f i"] & B
            \end{tikzcd}
            \quad \text{and} \quad
            \begin{tikzcd}
                X \arrow[d, equal] \arrow[r, "f"] & Y \arrow[d, "p"] \\
                X \arrow[r, "p f"] & B \arrow[d, equal] \\
                A \arrow[u] \arrow[r, "p f i"] & B
            \end{tikzcd}	
        \end{equation*}
        representing the same morphism $f \to p f i$ in $\TwAr(\sSet)$.\hfill$\Box$
        \medskip

        Let us now address the fibrancy properties of slices. Just as those of function spaces rested in the cofibrancy properties of the pushout-product operation, so do they rest on the cofibrancy properties of the pushout-join, which we recorded in \ref{chVI-3-18}, but did not yet prove. Here's the statement again:
        \medskip

        \ref{chVI-3-18} Lemma:
        Let $i \colon A \to B$ and $j \colon K \to L$ be cofibrations. Then $i \boxstar j$ is again a cofibration and
        \begin{enumerate}[label=\roman*)]
            \item[ii)-iii)]\addtocounter{enumi}{3}
                inner anodyne if either $i$ is right anodyne or $j$ is left anodyne
            \item
                left anodyne if $i$ is left anodyne
            \item
                right anodyne if $j$ is right anodyne
        \end{enumerate}
        \medskip

        We need one more lemma:

    \subsection{Lemma}\label{chVI-3-41}
        If $\Sigma \subseteq \mor (\sSet)$ is a saturated class of morphisms then
        \begin{equation*}
            \{ i \in \mor (\sSet) \mid i \boxstar j \in \Sigma \}
        \end{equation*}
        and
        \begin{equation*}
            \{ i \in \mor (\sSet)\ \mid j \boxstar i \in \Sigma \}
        \end{equation*}
        whenever $\Sigma$ is of the form
        \begin{equation*}
            \Sigma = \{ i \in \mor (\sSet) \mid i \text{ admits left lifts against } J \}
        \end{equation*}
        for some class $J \subseteq \mor (\sSet)$.
        \smallskip

        Note that \ref{chV-2-31} says that $\Sigma = \sat (S)$ for some set $S \subseteq \mor (\sSet)$ with compact sources satisfies the assumptions; in particular this is the case for $\Sigma$ the left/right/inner/$\varnothing$ anodyne maps.

    \subsection{Remark}\label{chVI-3-42}
        The analogue of this result for the pushout-product was exercise 3 on sheet 7. It was easier in that for every saturated class $\Sigma \in \mor (\sSet)$ the set
        \begin{equation*}
            \{ i \in \mor (\sSet) \mid i \boxstar j \in \Sigma \}
        \end{equation*}
        is again saturated, without the assumption that $\Sigma$ was defined by left lifting against some class $J$. The proof of that statement transfers to the present situation to show closure under retracts (trivially), pushouts and countable compositions, but not coproducts, since the join functor, when regarded as a functor $\sSet \to \sSet$ only preserves connected colimits.\\
        There are other sufficient conditions on $\Sigma$ such as right cancelability against cofibrations, i.e. if $i \in \Sigma$, $j$ cofibration and $i \circ j \in \Sigma$ implies $j \in \Sigma$. This is also true in all relevant examples , but nontrivial for all three cases of interest (though we know it for anodynes: By \ref{chV-4-4} and \ref{chV-2-27} we have
        \begin{equation*}
            \{ \text{anodyne extensions} \} = \{ \text{cofibrations }\} \cap \{ \text{weak homotopy equivalences}\}
        \end{equation*}	
        and the latter two are clearly right cancelable against all maps).

        \emph{Proof of \ref{chVI-3-41}}: By assumption, $i \boxstar j \in \Sigma$ is equivalent to admitting left lifts against $J$, i.e. every diagram
        \begin{equation*}
            \begin{tikzcd}
                A \star Y \underset{A \star X}{\cup} B \star X \arrow[d, "i \, \boxstar \, j"'] \arrow[r] & K \arrow[d, "f"] \\
                B \star Y \arrow[ur, dashed] \arrow[r] & L
            \end{tikzcd}
        \end{equation*}	
        admitting a filler for all $f \in J$. But by adjunction this is the same as
        \begin{equation*}
            \begin{tikzcd}
                A \arrow[d, "i"'] \arrow[r] & K_{/ B \to K} \arrow[d] \\
                B \arrow[ru, dashed] \arrow[r] & (K_{/ A \to K}) \underset{L_{/ A \to L}}{\times} (L_{/ B \to L})
            \end{tikzcd}
        \end{equation*}
        admitting fillers, which defines a saturated class by \ref{chV-2-2}.\\
        The other case is analogous using right instead of left slices.\hfill$\Box$

        \emph{Proof of \ref{chVI-3-18}}: Part i) is obvious (and we used it in \ref{chVI-3-30} already). For part ii) consider the class
        \begin{equation*}
            C = \{ k \in \mor (\sSet) \mid k \boxstar l \text{ is inner anodyne for all cofibrations } l\}
        \end{equation*}	
        By \ref{chVI-3-34} it is an intersection of saturated classes and thus itself saturated. We need to show that it contains the right anodynes, so it suffices to show
        \begin{equation*}
            \horn{n}{i} \to \Delta^n \in C
        \end{equation*}
        for all $0 < i \leq n$. But consider
        \begin{equation*}
            D = \{ l \in \mor (\sSet) \mid (\horn{n}{i} \to \Delta^n) \boxstar l \text{ is inner anodyne}\}
        \end{equation*}
        Again, this is saturated by \ref{chVI-3-34}, and we have to show that all cofibrations are contained in $D$, for which it suffices to treat the boundary inclusions $\partial \Delta^k \to \Delta^k$ by exercise 2, sheet 7. But lemma \ref{chVI-3-30} identifies $(\horn{n}{i} \to \Delta^n) \boxstar (\partial \Delta^k \to \Delta^k)$ as an inner horn, which is certainly inner anodyne.\\
        The argument in the other cases is the same.\hfill$\Box$

    \subsection{Proposition}\label{chVI-3-43}
        Let $A \xrightarrow{i} X \xrightarrow{f} Y \xrightarrow{p}$ be composable morphisms maps of simplicial sets, such that $i$ is a cofibration and $p$ as inner fibration, Then consider the maps
        \begin{equation*}
            s \colon Y_{/ f} \to Y_{/ f i} \underset{B_{/ p f i}}{\times} B_{/ p f}
        \end{equation*}
        and
        \begin{equation*}
            t \colon Y_{f /} \to Y_{f i /} \underset{B_{p f i /}}{\times} B_{p f /} \,. 
        \end{equation*}
        \begin{enumerate}[label=\roman*)]
            \item 
                $t$ is a left fibration and $s$ a right fibration
            \item 
                if $p$ is a right fibration, then $t$ is a Kan fibration, and if $p$ is a left fibration, then $s$ is a Kan fibration
            \item 
                if $p$ is a trivial fibration, then so are $t$ and $s$
            \item 
                if $i$ is right anodyne, then $t$ is a trivial fibration, and if $i$ is left anodyne then $s$ is a trivial fibration.
        \end{enumerate}

        \emph{Proof}: We prove statements about $t$: As discussed, a lifting problem
        \begin{equation*}
            \begin{tikzcd}
                S \arrow[d, "a"'] \arrow[r] & Y_{f /} \arrow[d, "t"] \\
                T \arrow[ru, dashed] \arrow[r] & Y_{f i /} \underset{B_{p f i /}}{\times} B_{p f /}
            \end{tikzcd}	
        \end{equation*}
        is equivalent to
        \begin{equation*}
            \begin{tikzcd}
                A \star T \underset{A \star S}{\cup} X \star S \arrow[d, "i \, \boxstar \, a"'] \arrow[r] & Y \arrow[d, "p"] \\
                X \star T \arrow[ru, dashed] \arrow[r] & B
            \end{tikzcd}	
        \end{equation*}
        \begin{enumerate}[label=\roman*)]
            \item 
                To show that $t$ is a left fibration we consider the inclusions $\horn{n}{i} \to \Delta^n$ for $0 \leq i < n$ for $a$, which are left anodyne by definition. Thus \ref{chVI-3-18} shows that $i \boxstar a$ is then inner anodyne, so since $p$ is an inner fibration, the lower diagram admits a lift.
            \item 
                If we take $a = \horn{n}{n} \to \Delta^n$ then \ref{chVI-3-18} gives that $i \boxstar a$ is still right anodyne so a lift exists whenever $p$ is a right fibration.
            \item 
                If we take $a = \partial \Delta^n \to \Delta^n$ then $i \boxstar a$ is still a cofibration, so a lift exists if $p$ is a trivial fibration.
            \item 
                If again $a = \partial \Delta^n \to \Delta^n$, but $i$ is right anodyne, then $i \boxstar a$ is inner anodyne by \ref{chVI-3-18}, so a lift always exists.\hfill$\Box$
        \end{enumerate}

    \subsection{Corollary}\label{chVI-3-44}
        If $f \colon \Ii \to \Cc$ is a map to a quasicategory then
        \begin{equation*}
            s \colon \Cc_{/ f} \to \Cc \qquad t \colon \Cc_{f /} \to \Cc
        \end{equation*}	
        are a right and left fibration, respectively.\\
        In particular, $\Cc_{/ f}$ and $\Cc_{f /}$ are again quasicategories.
        \medskip

        Before we finish the proof of Joyal's theorem, a remark:

    \subsection{Remark}\label{chVI-3-45}
        We now have two competing definitions for over- and under-categories of quasicategories: On the one hand for $X \colon \Delta^0 \to \Cc$ there are the categories
        \begin{equation*}
            \Cc_{/X} \quad \text{and} \quad \Cc_{X /}
        \end{equation*}
        defined using joins and on the other hand there are
        \begin{equation*}
            \Cc \big/ X \quad \text{and} \quad X \big/ \Cc
        \end{equation*}
        defined using the arrow category. Worse yet, the latter construction can be extended as follows: Given $F \colon \Ii \to \Cc$ and $G \Jj \to \Cc$ put
        \begin{equation*}
            \begin{tikzcd}
                F \big/ G \arrow[d] \arrow[r] \arrow[rd, phantom, "\lrcorner", very near end] & \Ar(\Cc) \arrow[d, "(s{,}t)"] \\
                \Ii \times \Jj \arrow[r] & \Cc \times \Cc
            \end{tikzcd}	
        \end{equation*}
        As $(s, t)$ is an inner fibration $f \big/ G$ is a quasicategory whenever $\Ii$ and $\Jj$ are such, and it follows from \ref{chVI-3-31} that the map $f \big/ G \to \Ii \times \Jj$ is a left fibration if $\Ii \simeq *$ and a right fibration if $\Jj \simeq *$. But then 
        \begin{equation*}
            \Cc_{\sslash f} \coloneq \const \colon \Cc \to \Fun(\Ii, \Cc) \Big/ f \colon \Delta^0 \to \Fun(\Ii, \Cc)
        \end{equation*}
        is another perfectly reasonable attempt at defining a category of cones over a functor $f \colon \Ii \to \Cc$.\\
        It will be  an exercise to show, that $\Cc_{\sslash f}$ and $\Cc_{/ f}$ agree for nerves of ordinary categories, but this is not true in general. Furthermore, since by definition both
        \begin{equation*}
            \begin{tikzcd}
                \Hom_{\Cc}(Y, X) \arrow[d] \arrow[r] & \Cc \big/ X \arrow[d] \\
                \Delta^0 \arrow[r, "Y"] & \Cc
            \end{tikzcd}
            \quad
            \text{and}
            \quad
            \begin{tikzcd}
                \Hom_{\Cc}(Y, X) \arrow[d] \arrow[r] & Y \big/ \Cc \arrow[d] \\
                \Delta^0 \arrow[r, "X"] & \Cc
            \end{tikzcd}	
        \end{equation*}
        are pullbacks, the diagrams
        \begin{equation*}
            \begin{tikzcd}
                \Hom_{\Cc}^R(Y, X) \arrow[d] \arrow[r] \arrow[rd, phantom, "\lrcorner", very near end] & \Cc_{/ X} \arrow[d] \\
                \Delta^0 \arrow[r, "Y"] & \Cc
            \end{tikzcd}
            \quad
            \text{and}
            \quad
            \begin{tikzcd}
                \Hom_{\Cc}^L(Y, X) \arrow[d] \arrow[r] \arrow[rd, phantom, "\lrcorner", very near end] & \Cc_{Y /} \arrow[d] \\
                \Delta^0 \arrow[r, "X"] & \Cc
            \end{tikzcd}	
        \end{equation*}
        give reasonable definitions of mapping spaces (they are Kan by \ref{chVI-3-44} and \ref{chVI-3-27}). As the notation suggests these are different from $\Hom_{\Cc}(Y, X)$ and inn fact from one another! There are natural injections
        \begin{equation*}
            \Hom_{\Cc}^R(Y, X) \hookrightarrow \Hom_{\Cc}(Y, X) \hookleftarrow \Hom_{\Cc}^L(Y, X)
        \end{equation*}
        but it is a non-trivial result that these are homotopy equivalences, and similarly, that the two categories of cones are in fact equivalent.\\
        The formation of fat slice categories also admits a left adjoint, the \emph{fat join}, so one may wonder, why bother with the usual join and slices at all... The answer is that necessarily the fat join of $\Delta^n$ and $\Delta^m$ is \emph{not} $\Delta^{n+m+1}$, so the one thing they cannot obviously be used for is a proof of Joyal's lifting theorem! Quite the pickle...

        These issues will come facts to haunt us when establishing Theorem H, as $\Hom_{\Cc}(X, Y)$ is more easily compared to $\Hom_{\Nn^c(\Cc)}^L(X, Y)$ than $\Hom_{\Nn^c(\Cc)}(X, Y)$.
        \medskip

        We are finally ready to prove Joyal's theorem:

    \subsection{Proof of \ref{chvI-3-20} (existence)}
        By \ref{chVI-3-35} the lifting problem
        \begin{equation*}
            \begin{tikzcd}
                \horn{n}{0} \arrow[d] \arrow[r, "f"] & \Cc \arrow[d, "F"] \\
                \Delta^n \arrow[ru, dashed] \arrow[r] & \Dd
            \end{tikzcd}	
        \end{equation*}	
        is equivalent to
        \begin{equation*}
            \begin{tikzcd}
                \Delta^0 \arrow[d, "d_1"'] \arrow[r] & \Cc_{/ \Delta^{n-2}} \arrow[d] \\
                \Delta^1 \arrow[ru, dashed] \arrow[r] & \Cc_{/ \partial \Delta^{n-2}} \underset{\Dd_{\partial \Delta^{n-2}}}{\times} \Dd_{/ \Delta^{n-2}}
            \end{tikzcd}	
        \end{equation*}
        As $F$ is an inner fibration and $\partial \Delta^{n-2} \to \Delta^{n-2}$ a cofibration the right hand vertical map is a right fibration by \ref{chVI-3-43}. Furthermore, by construction there is a commutative diagram
        \begin{equation*}
            \begin{tikzcd}
                \Delta^1 \arrow[dd, "(01)"'] \arrow[r] & \Cc_{/ \partial \Delta^{n-2}} \underset{\Dd_{\partial \Delta^{n-2}}}{\times} \Dd_{/ \Delta^{n-2}} \arrow[d] \arrow[r] \arrow[rd, phantom, "\ulcorner" very near end] & \Dd_{/ \Delta^{n-2}} \arrow[d, color=magenta] \\
                & \Cc_{/ \partial \Delta^{n-2}} \arrow[d, color=magenta] \arrow[r] & \Dd_{/ \partial \Delta^{n-2}} \\
                \horn{n}{0} \arrow[r, "f"] & \Cc &
            \end{tikzcd}
        \end{equation*}
        By \ref{chVI-3-43} these two magenta vertical maps are right fibrations, thus so is the middle vertical composite. This implies first that $\Cc_{/ \partial \Delta^{n-2}} \underset{\Dd_{\partial \Delta^{n-2}}}{\times} \Dd_{/ \Delta^{n-2}}$ is a quasicategory and secondly \ref{chV-38} implies that
        \begin{equation*}
            \Delta^1 \to \Cc_{/ \partial \Delta^{n-2}} \underset{\Dd_{\partial \Delta^{n-2}}}{\times} \Dd_{/ \Delta^{n-2}}
        \end{equation*}
        represents an equivalence as the composite
        \begin{equation*}
            \Delta^1 \to \horn{n}{0} \xrightarrow{f} \Cc
        \end{equation*}
        does so by assumption. But then \ref{chVI-3-36} and \ref{chVI-3-38} give the desired lift.\hfill$\Box$
        \medskip

        To establish uniqueness of the lift in Joyal's theorem it suffices to show:

    \subsection{Lemma}\label{chVI-3-46}
        If $p \colon \Cc \to \Dd$ is a right fibration between quasicategories, then the fibre of the map
        \begin{equation*}
            r \,\colon \Fun(\Delta^1, \Cc) \to \Fun(\Delta^1, \Dd)_s \underset{\Dd}{\times} \Cc
        \end{equation*}
        over a pair $(f, x)$ is a contractible Kan complex, whenever $f \colon \Delta^1 \to \Dd$ is an equivalence.
        \medskip

        The proof will use \ref{chVI-3-27}, which in turn relies only on the existence part of \ref{chVI-3-20}, so we avoid circularity.
        \smallskip

        \emph{Proof}: Since $p$ is conservative by \ref{chVI-3-28}, the diagram
        \begin{equation*}
            \begin{tikzcd}
                \core(\Cc) \arrow[d, "\core(p)"'] \arrow[r] & \Cc \arrow[d, "p"] \\
                \core(\Dd) \arrow[r] & \Dd
            \end{tikzcd}	
        \end{equation*}
        is a pullback. In particular, $\core(p)$ is again a right fibration, and thus a Kan fibration by \ref{chVI-3-27}. But then by \ref{chV-2-26} the map
        \begin{equation*}
            \overline{r} \,\colon \Fun(\Delta^1, \core(\Cc)) \to \Fun(\Delta^1, \core(\Dd)) \underset{\Fun(\{0\}, \core(\Dd))}{\times} \Fun(\{0\}, \core(\Cc))
        \end{equation*}
        is a trivial fibration, so in particular has contractible fibres. It remains to check that
        \begin{equation*}
            \overline{r}^{-1}(f, x) = r^{-1}(f, x)
        \end{equation*}
        whenever $f$ is an equivalence. But these simplicial sets have the same vertices and edges as $p$ is conservative. This is immediate for the vertices and for edges one decodes that an edge of $r^{-1}(f, x)$ is a diagram $\Delta^1 \times \Delta^1 \to \Cc$:
        \begin{equation*}
            \begin{tikzpicture}[
                scale=0.9, 
                thick, 
                color=violet, 
                every node/.style={text=black}
            ]
                    \coordinate (A) at (0, 0);
                    \coordinate (B) at (0, 2);
                    \coordinate (C) at (2, 2);
                    \coordinate (D) at (2, 0);

                    \node[below] at (A) {$x$};
                    \node[below] at (D) {$x$};                
                    \draw[postaction={nomorepostaction,decorate,decoration={markings,mark=at position 0.6 with {\arrow{>[length=8pt,width=5pt]}}}}] (A) -- node[left] {$g$} (B);
                    \draw[postaction={nomorepostaction,decorate,decoration={markings,mark=at position 0.6 with {\arrow{>[length=8pt,width=5pt]}}}}] (B) -- node[above] {$s$} (C);
                    
                    \draw[postaction={nomorepostaction,decorate,decoration={markings,mark=at position 0.6 with {\arrow{>[length=8pt,width=5pt]}}}}] (A) -- (C);

                    \draw[postaction={nomorepostaction,decorate,decoration={markings,mark=at position 0.6 with {\arrow{>[length=8pt,width=5pt]}}}}] (A) -- node[below] {$\id_x$} (D);

                    \draw[postaction={nomorepostaction,decorate,decoration={markings,mark=at position 0.6 with {\arrow{>[length=8pt,width=5pt]}}}}] (D) -- node[right] {$g'$} (C);

                    %\draw[fill=violet,fill opacity=0.3] (A) -- (B) -- (C) -- cycle;
                    %\draw[fill=violet,fill opacity=0.3] (A) -- (C) -- (D) -- cycle;
                    
            \end{tikzpicture}
        \end{equation*}
        where $g, g'$ are lifts of $f$ and $s$ lifts the identity of $d_0(f)$. In particular all these edges are equivalences by conservativity of $p$, so the entire diagram factors through $\core(\Cc)$. For higher simplices $\Delta^n \to r^{-1}(f, x)$, note that the condition that
        \begin{equation*}
            \Delta^1 \times \Delta^1 \to \Cc
        \end{equation*}
        factors through $\core(\Cc)$ depends only on the images of edges, which we just checked.\hfill$\Box$

    \subsection{Proof of \ref{chVI-3-20} (uniqueness)}
        Apply \ref{chVI-3-46} to the reformulation of the lifting problem in the proof of existence.\hfill$\Box$
        \medskip

        Finally, we deduce the formulation of Joyal's theorem given in \ref{chVI-3-16}. Given a morphism $f \colon x \to y$ in a quasicategory $\Cc$ consider the functors
        \begin{equation*}
            \Cc_{/ x} \leftarrow \Cc_{/ f} \rightarrow \Cc_{/ y} 
        \end{equation*}
        arising from the sequences $\Delta^0 \xrightarrow{0 \text{ or } 1} \Delta^1 \xrightarrow{f} \Cc \rightarrow *$. Then by \ref{chVI-3-43} the left functor is a trivial Kan fibration, so in particular an equivalence of categories by \ref{chVI-2-12}. Choosing an inverse (of which there is a contractible choice by \ref{chVI-3-24}) we obtain a functor
        \begin{equation*}
            f_* \,\colon \Cc_{/ x} \to \Cc_{/ y}
        \end{equation*}
        that on vertices is given by postcomposition with $f$. Analogously, there is a functor
        \begin{equation*}
            f^* \,\colon \Cc_{y /} \to \Cc_{x /}
        \end{equation*}
        given by precomposition with $f$. We then have the following refinement of \ref{chVI-3-20}:

    \subsection{Proposition}\label{chVI-3-47}
        For a morphism $f \colon x \to y$ in a quasicategory $\Cc$ the statements
        \begin{enumerate}[label=\roman*)]
            \item 
                $f$ is an equivalence \label{chVI-3-47-i}
            \item 
                $\Cc_{/ f} \to \Cc_{/ y}$ is a trivial Kan fibration \label{chVI-3-47-ii}
            \item 
                $\Cc_{/ f} \to \Cc_{/ y}$ is an equivalence of categories \label{chVI-3-47-iii}
            \item 
                'the' functor $f_* \colon \Cc_{/ x} \to \Cc_{/ y}$ is an equivalence of categories \label{chVI-3-47-iv}
        \end{enumerate}	
        are equivalent.
        \smallskip

        The analogous statements using under-categories hold as well.
        \smallskip

        \emph{Proof}: $\ref{chVI-3-47-i} \to \ref{chVI-3-47-ii}$: A lifting problem
        \begin{equation*}
            \begin{tikzcd}
                \partial \Delta^n \arrow[d] \arrow[r] & \Cc_{/ f} \arrow[d] \\
                \Delta^n \arrow[r] & \Cc_{/ x}
            \end{tikzcd}	
        \end{equation*}
        by using \ref{chVI-3-35} translates to
        \begin{equation*}
            \begin{tikzcd}
                \Delta^1 \arrow[r] \arrow[rr, "f", bend left] & \horn{n+2}{0} \arrow[d] \arrow[r] & \Cc \arrow[d] \\
                & \Delta^{n+2} \arrow[r] & * 
            \end{tikzcd}	
        \end{equation*}
        which has a solution by \ref{chVI-3-20} if $f$ is an equivalence. \\
        $\ref{chVI-3-47-ii} \to \ref{chVI-3-47-iii}$: \ref{chVI-2-12}. \\
        $\ref{chVI-3-47-iii} \to \ref{chVI-3-47-iv}$: Immediate from the homotopy commutative diagram
        \begin{equation*}
            \begin{tikzcd}
                \Cc_{/ x} \arrow[rd, "f_*"] & \Cc_{/ f} \arrow[l] \arrow[d] \\
                & \Cc_{/ y}
            \end{tikzcd}	
        \end{equation*}	
        since the horizontal map is an equivalence.
        $\ref{chVI-3-47-iv} \to \ref{chVI-3-47-i}$: If $f_* \colon \Cc_{/ x} \to \Cc_{/ y}$ is an equivalence of categories we can in particular pick a $g \in \Cc_{/ x}$ such that $f_* g$ is equivalent to $\id_y$ in $\Cc_{/ y}$. But decoding this means that there is a $2$-simplex in $\Cc$
        \begin{equation*}
            \begin{tikzpicture}[
                scale=0.9, 
                thick, 
                color=violet, 
                every node/.style={text=black}
            ]
                    \coordinate (A) at (0, 0);
                    \coordinate (B) at (1, 1.73);
                    \coordinate (C) at (2, 0);

                    \node[below] at (A) {$y$};
                    \node[below] at (C) {$y$};                
                    \draw[postaction={nomorepostaction,decorate,decoration={markings,mark=at position 0.6 with {\arrow{>[length=8pt,width=5pt]}}}}] (A) -- node[left] {$s$} (B);
                    \draw[postaction={nomorepostaction,decorate,decoration={markings,mark=at position 0.6 with {\arrow{>[length=8pt,width=5pt]}}}}] (B) -- node[right] {$h$} (C);
                    \draw[postaction={nomorepostaction,decorate,decoration={markings,mark=at position 0.6 with {\arrow{>[length=8pt,width=5pt]}}}}] (A) -- node[below] {$\id_y$} (C);
                    
            \end{tikzpicture}        
        \end{equation*}	
        where $s$ is an equivalence and $h$ a composite of $f$ and $g$. But this means that $\id_y$ is a composite of $f$ and $g s$. So $f$ admits a right inverse. But then the images of $\id_x$ and $g s f$ in $\Cc_{/ y}$ are equivalent since
        \begin{equation*}
            f_* \id_x \sim f \sim \id_y f \sim f g s f \sim f_* g s f \,.
        \end{equation*}
        But as $f_*$ is an equivalence this implies that $\id_y$ and $g s f$ are equivalent in $\Cc_{/ x}$. Decoding again, we find this implies that $g s f$ is an equivalence, so $f$ also admits a left inverse, hence is an equivalence.\hfill$\Box$

    \subsection{Remark}\label{chVI-3-48}
        We have \emph{not} constructed a functor
        \begin{equation*}
            \Cat_{\infty \Cc_{/}} \to \Cat_{\infty}, \quad f \colon \Cc \to \Dd \mapsto \Dd_{f /} \text{ or } \Dd_{/ f}
        \end{equation*}
        or anything similar so far! This is quite difficult, in fact already showing that an equivalence of categories $g \colon \Dd \to \Dd'$ gives equivalences
        \begin{equation*}
            \Dd_{f /} \to \Dd'_{g f /} \quad \text{and} \quad \Dd_{/ f} \to \Dd'_{/ g f}
        \end{equation*}
        is a highly non-trivial endeavor. Here again the fat slices are better behaved; the latter statement is almost a triviality here. To obtain a functor one 'only' needs to understand the interaction of the coherent nerve functor with over- and under-categories to build a simplicial model for $\Cat_{\infty \Cc_{/}}$; alternatively the statement will follow from general machinery about limits in quasicategories.

        In fact the statement about thin slices is usually proved by reduction to the corresponding one for fat slices!
